\section{Przebieg gry}

Gra ,,Labirynt śmierci'' polega na wykonywaniu przez graczy określonych czynności w ściśle określonej kolejności.
Jej celem jest odnalezienie przez oddział i zniszczenie magicznych Czarnych Wrót.

Decyzje, dotyczące zachowania się postaci podczas gry mogą być podejmowane wspólnie przez wszystkich graczy lub też, jeżeli tak zostanie zadecydowane, każdy z graczy może podejmować decyzje indywidualne.

,,Labirynt śmierci'' rozgrywany jest w etapach.
Każdy etap jest to blok czynności, które gracze powinni wykonać przed przystąpieniem do następnego etapu.
Czynności te muszą być wykonywane w przypisanej im niżej kolejności:

\subsection{Umieszczenie nowego segmentu}
\begin{enumerate}
    \item
        Oddział ustala drogę, którą opuści zajmowany obecnie segment, czyli kierunek, w którym pójdzie.
    \item
        Jeden z graczy losowo wybiera z pojemnika nowy segment labiryntu (żeton).
    \item
        Nowy segment układany jest obok tego, który oddział opuszcza, na przedłużeniu drogi, którą zdecydował się on wybrać.
\end{enumerate}

Nowy segment musi zostać ułożony tak, by pasował do segmentu opuszczanego, tzn. korytarz musi stanowić przedłużenie korytarza, a przejście musi łączyć się z przejściem.
Od decyzji graczy zależy w jaki sposób zostanie ułożony taki segment, w przypadku którego możliwe jest więcej niż jedno prawidłowe ułożenie.

\subsection{Poszukiwanie pułapek}
Poszukiwanie pułapek (rzucanie kostką) odbywa się tylko wtedy, gdy oddział chce wejść do nieodwiedzanej dotychczas komnaty, bądź też opuszcza komnatę inną drogą, niż do niej wszedł.
A więc oddział nie szuka pułapek wtedy, gdy porusza się korytarzem.
\begin{enumerate}
    \item
        Jeden z graczy rzuca 1K6, by sprawdzić, czy drzwi, przez które oddział ma zamiar przejść, zawierają pułapkę.
        Drzwi zawierają pułapkę, jeżeli wyrzucone zostanie 1 oczko.
    \item
        Po odnalezieniu pułapki (wyrzucone 1 oczko) należy wydelegować jedną postać do jej zbadania.
        Delegowana postać powinna posiadać w grupie zdolności hasło ,,pułapki'', gdyż tylko wtedy może podjąć próbę unieszkodliwienia znalezionej pułapki (patrz 7.0).
    \item
        Jeżeli próba unieszkodliwienia pułapki zakończy się niepowodzeniem lub w ogóle nie zostanie podjęta uważa się, że pułapka zadziałała.
        Gracze muszą ustalić jej rodzaj, rzucając 1K6 i odczytując rezultat z tabeli 7.1 ,,Pułapki''.
\end{enumerate}

\subsection{Spotkanie z potworem}
Oddział wykonuje ruch.
Odbywa się to przez umieszczenie jego żetonu na nowym segmencie.
Następną czynnością jest sprawdzenie, czy w nowym segmencie znajdują się potwory.
Jeden z graczy rzuca 1K6.
Spotkanie z potworem następuje przy wyrzuceniu:

1, 2 lub 3 oczek jeżeli nowy segment jest nie zbadaną dotychczas komnatą,

1 oczka jeżeli nowy segment jest fragmentem korytarza lub komnatą, którą oddział już poprzednio odwiedził.

Oddział nie musi od razu przystępować do walki z napotkanym potworem.
Może przedtem podjąć próbę negocjacji lub wykupienia się (patrz 8.0).

Jeżeli próby negocjacji lub wykupienia się zawiodą lub nie zostaną podjęte rozpoczyna się walka.
Składa się ona z następujących faz (czynności):
\begin{enumerate}
    \item
        atak oddziału,
    \item
        atak potwora,
    \item
        atak Czarnych Wrót,
    \item
        reorganizacja szyku oddziału,
    \item
        reorganizacja szyku grupy potworów.
\end{enumerate}

Powyższe fazy muszą być wykonywane w podanej kolejności.

Po zabiciu potwora (lub wszystkich potworów, stanowiących napotkaną grupę) oddział ma prawo odebrać mu posiadane przez niego skarby.
Ustala się ile i jakie skarby potwór posiadał (patrz 14.0), a następnie dzieli się je między członków oddziału.

Jeżeli potwór zostanie zabity pozostali przy życiu członkowie oddziału zdobywają określoną liczbę punktów doświadczenia (patrz 12.0).

\subsection{Badanie znalezisk}

Znaleziska są to rzeczy, znalezione przez oddział w komnatach (tylko).

\begin{enumerate}
    \item
        Jedna z postaci jest delegowana do zbadania znaleziska.
    \item
        Badająca postać (gracz) rzuca 1K6.
        Z tabeli 13.9 ,,Znaleziska'' z kolumny, odpowiadającej jego rodzajowi, odczytywana jest dokładna nazwa tego, co zostało znalezione.
    \item
        Wprowadzane są w życie wszelkie efekty, wiążące się z danym znaleziskiem.
\end{enumerate}

Wszystkie wymienione powyżej czynności nazywane są wspólnie etapem gry.
Po ich wykonaniu gracze przystępują do kolejnego etapu, czyli od początku do umieszczenia kolejnego segmentu labiryntu.
